\documentclass[11pt]{article}

\usepackage[T1]{fontenc}
\usepackage[utf8]{inputenc}
\usepackage{amsmath,amssymb,amsthm,mathtools}
\usepackage[colorlinks=true,linkcolor=blue,citecolor=blue,urlcolor=blue]{hyperref}

\newtheorem{theorem}{Theorem}
\newtheorem{proposition}{Proposition}
\newtheorem{lemma}{Lemma}
\newtheorem{corollary}{Corollary}
\theoremstyle{definition}
\newtheorem{definition}{Definition}
\newtheorem{problem}{Problem}
\theoremstyle{remark}
\newtheorem{remark}{Remark}

\DeclareMathOperator{\Var}{Var}
\DeclareMathOperator{\Prob}{P}
\newcommand{\E}{\mathbb E}
\newcommand{\N}{\mathbb N}
\newcommand{\Z}{\mathbb Z}
\newcommand{\1}{\mathbf 1}

\title{Zeta-Law Entropy and Zeta-Inequality Framework}
\author{Pudim AI Project}
\date{}

\begin{document}
\maketitle

\begin{abstract}
This note records the audited core of a zeta-law framework for treating zeta-ratio identities, modular entropy shadows, and applications to inequalities for the Riemann zeta function.  The fully verified material in this version consists of the Gibbs-law formulation, free-energy notation, modular and Mellin-Planck definitions, and Euler-score identities.  Several stronger application claims are recorded as audit targets.
\end{abstract}

\section{The zeta law}

\begin{definition}[Riemann zeta probability law]
For \(\beta>1\), define the probability law \(\rho_\beta\) on \(\N\) by
\[
\rho_\beta(n)=\frac{n^{-\beta}}{\zeta(\beta)},\qquad n\in \N.
\]
Equivalently,
\[
\rho_\beta(n)=\frac{e^{-\beta E(n)}}{Z(\beta)},\qquad
E(n)=\log n,\qquad Z(\beta)=\zeta(\beta).
\]
\end{definition}

\begin{definition}[Zeta free energy]
The zeta free energy is
\[
A(\beta)=\log \zeta(\beta),\qquad \beta>1.
\]
Differentiating under the absolutely convergent Dirichlet series gives
\[
A'(\beta)=-\E_\beta[\log N],
\qquad
A''(\beta)=\Var_\beta(\log N).
\]
\end{definition}

\begin{remark}
For \(r\ge 0\), the zeta law satisfies the moment identity
\[
\E_\beta[N^{-r}]=\frac{\zeta(\beta+r)}{\zeta(\beta)}.
\]
Thus many zeta-ratio inequalities can be read as moment inequalities under \(\rho_\beta\).
\end{remark}

\section{Modular successor entropy}

\begin{definition}[Residue distribution and successor entropy]
For \(q\ge2\), define the residue distribution of the zeta law modulo \(q\) by
\[
\mu_{q,\beta}(a)=
\sum_{\substack{n\ge1\\ n\equiv a\pmod q}}\rho_\beta(n),
\qquad a\in \Z/q\Z.
\]
Writing \(a+1\) cyclically in \(\Z/q\Z\), define
\[
B_q(\beta)=
\sum_{a\in \Z/q\Z}
\mu_{q,\beta}(a)
\log\frac{\mu_{q,\beta}(a)}{\mu_{q,\beta}(a+1)}.
\]
\end{definition}

\begin{problem}[Successor entropy resolution]
Audit the claimed identity
\[
\Sigma(\beta)
=\sum_{n=1}^{\infty}\rho_\beta(n)
\log\frac{\rho_\beta(n)}{\rho_\beta(n+1)}
=\frac{\beta}{\zeta(\beta)}
\sum_{k=1}^{\infty}\frac{(-1)^{k+1}}{k}\zeta(\beta+k),
\]
together with the modular variational formula
\[
\Sigma(\beta)=\sup_{q\ge2}B_q(\beta)=\lim_{q\to\infty}B_q(\beta).
\]
The intended proof expands \(\log(1+1/n)\), uses the log-sum inequality, and lets large moduli isolate finite initial intervals.
\end{problem}

\begin{problem}[Prime-modulus Dirichlet \(L\)-resolution]
For prime \(p\), let \(\chi\) range over Dirichlet characters modulo \(p\), and set
\[
S_a(p,\beta)=\sum_{\chi\bmod p}\chi(a)L(\beta,\chi).
\]
Audit the claimed resolution
\[
\mu_{p,\beta}(a)=\frac{S_a(p,\beta)}{(p-1)\zeta(\beta)}
\quad (1\le a\le p-1),
\qquad
\mu_{p,\beta}(0)=p^{-\beta},
\]
and the resulting prime-modulus formula for the successor entropy.
\end{problem}

\section{Euler-score identities}

\begin{proposition}[Euler-score decomposition]
For every \(\beta>1\),
\[
-\bigl(\log\zeta\bigr)'(\beta)
=\E_\beta[\log N]
=\sum_{d=1}^{\infty}\Lambda(d)\mu_{d,\beta}(0)
=\sum_{d=1}^{\infty}\frac{\Lambda(d)}{d^\beta}.
\]
Equivalently,
\[
-\frac{\zeta'(\beta)}{\zeta(\beta)}
=\sum_{d=1}^{\infty}\frac{\Lambda(d)}{d^\beta}.
\]
\end{proposition}

\begin{proof}
Use the von Mangoldt identity
\[
\log n=\sum_{d\mid n}\Lambda(d).
\]
Then
\[
\E_\beta[\log N]
=\sum_{n\ge1}\rho_\beta(n)\sum_{d\mid n}\Lambda(d)
=\sum_{d\ge1}\Lambda(d)\Prob_\beta(d\mid N).
\]
For the zeta law,
\[
\Prob_\beta(d\mid N)
=\frac{1}{\zeta(\beta)}\sum_{m\ge1}(dm)^{-\beta}
=d^{-\beta}.
\]
The identity with \(-(\log\zeta)'(\beta)\) follows from the free-energy derivative.
\end{proof}

\begin{proposition}[Finite Euler-score identity]
For every \(A\subseteq\{1,\ldots,N\}\),
\[
\sum_{n\in A}\log n
=\sum_{d\le N}\Lambda(d)\sum_{m\le N/d}\1_A(md).
\]
\end{proposition}

\begin{proof}
Reverse the order of summation:
\[
\sum_{d\le N}\Lambda(d)\sum_{m\le N/d}\1_A(md)
=\sum_{n\in A}\sum_{d\mid n}\Lambda(d)
=\sum_{n\in A}\log n.
\]
\end{proof}

\section{Mellin-Planck partition function}

\begin{definition}[Mellin-Planck partition function]
For \(s>1\), define
\[
M(s)=\Gamma(s)\zeta(s)
=\int_0^\infty \frac{t^{s-1}}{e^t-1}\,dt.
\]
The normalized law associated with \(M\) is
\[
\nu_s(dt)=\frac{t^{s-1}}{(e^t-1)M(s)}\,dt.
\]
\end{definition}

\begin{remark}
The free-energy identity
\[
\frac{d^2}{ds^2}\log M(s)=\Var_{\nu_s}(\log t)\ge 0
\]
shows that \(M\) is log-convex on \((1,\infty)\).
\end{remark}

\section{Application audit targets}

\begin{problem}[Nantomah positivity]
Audit the claimed proof that, for every \(n\in\N\),
\[
(n+2)\zeta(n+1)\zeta(n+3)
-(n+1)\zeta(n+2)^2
-\zeta(n+1)\zeta(n+2)>0.
\]
The imported route rewrites the expression using moments of \(X(m)=1/m\) under \(\rho_s\), then splits the result into a positive linear part and a bounded quadratic correction.
\end{problem}

\begin{problem}[Alzer-Kwong reciprocal-zeta curvature]
Let \(F(x)=1/\zeta(x)\).  Audit the claimed sign pattern
\[
F''(x)>0\quad\text{on }(-4n,-4n+2),
\qquad
F''(x)<0\quad\text{on }(-4n-2,-4n)
\]
for every integer \(n\ge1\).  The imported route uses the zeta functional equation and a positive-axis curvature bound for \((\log\zeta)''(s)\).
\end{problem}

\begin{problem}[Sroysang generalized Holder inequality]
Let \(m\ge2\), \(r\ge1\), \(x_1,\ldots,x_m>1\), and \(p_1,\ldots,p_m>1\) satisfy
\[
\sum_{i=1}^m\frac1{p_i}=\frac1r.
\]
With
\[
T=r\sum_{i=1}^m\frac{x_i}{p_i},
\]
audit the claimed inequality
\[
\zeta(T)\le
\frac{\prod_{i=1}^m(\Gamma(x_i)\zeta(x_i))^{r/p_i}}{\Gamma(T)}.
\]
The imported route applies generalized Holder to
\[
f_i(t)=\left(\frac{t^{x_i-1}}{e^t-1}\right)^{1/p_i}.
\]
\end{problem}

\section{Program}

The framework currently separates four proof layers:
\[
\frac{\zeta(s+r)}{\zeta(s)}=\E_s[N^{-r}],
\qquad
(\log\zeta)''(s)=\Var_s(\log N),
\]
\[
M(s)=\Gamma(s)\zeta(s),
\qquad
(\log M)''(s)=\Var_{\nu_s}(\log t),
\]
and
\[
\Sigma(\beta)=\sup_{q\ge2}B_q(\beta).
\]
The next mathematical task is to audit the imported application proofs and promote each application from problem status to theorem status only after the proof is self-contained and all endpoint, exponent, and extraction-sensitive formulas have been checked.

\end{document}
